% Texign worksheet example. Build with: texign build fisa.tex
% answers=hide (student copy, default), show (teacher copy), end (key at the end)
\documentclass{texign-worksheet}

\subject{Istorie modernă · Fișă de lucru}
\title{Școala românească după 1864}
\date{}
\author{}

\begin{document}
\maketitle

\begin{instructions}
Citiți sursa, apoi rezolvați exercițiile. Timp de lucru: 40 de minute.
\end{instructions}

\begin{source}{Legea instrucțiunii publice, 1864, art.~31}
„Instrucțiunea primară este obligatorie și gratuită.”
\end{source}

\begin{exercise}[title=Înțelegerea sursei]
\begin{questions}
\question[2] Ce principii stabilește articolul citat?
\answerlines[2]
\begin{answer}
Obligativitatea și gratuitatea învățământului primar.
\end{answer}
\question[3] Ce nu precizează textul legii despre aplicarea ei?
\answerlines[3]
\begin{answer}
Cine plătește învățătorul, cine construiește școala, cine verifică frecvența.
\end{answer}
\end{questions}
\end{exercise}

\begin{exercise}[title=Alegere multiplă, points=1]
În ce an a fost adoptată Legea instrucțiunii publice?
\begin{choices}
\choice 1848
\correct 1864
\choice 1896
\choice 1918
\end{choices}
\end{exercise}

\begin{exercise}[title=Completare, points=2]
Legea din \blank[20mm]{1864} a introdus învățământul primar obligatoriu; reformele ulterioare sunt legate de numele lui \blank[35mm]{Spiru Haret}.
\end{exercise}

\begin{exercise}[title=Argumentare, points=20]
Explicați, în 8--10 rânduri, de ce o lege nu descrie întotdeauna realitatea.
\answerbox[45mm]
\begin{answer}
Se acceptă orice răspuns care distinge norma de aplicarea ei.
\end{answer}
\end{exercise}

\begin{exercise}[title=Schiță, points=2]
Reprezentați grafic evoluția numărului de școli între 1864 și 1914.
\answergrid[40mm]
\end{exercise}

\begin{activity}[Discuție în perechi]
Comparați răspunsurile de la exercițiul~4 și formulați o concluzie comună.
\end{activity}
\end{document}
